\documentclass[a4paper,12pt]{article}

\usepackage[utf8]{inputenc}
\usepackage[T1]{fontenc}
\usepackage[italian]{babel}
\usepackage{subcaption}

\usepackage{booktabs}
\usepackage{geometry}

\geometry{margin=2.0cm}
\usepackage{setspace}
\setstretch{1.3}

\usepackage{amsmath,amssymb}
\usepackage{siunitx}
\usepackage{graphicx}
\usepackage{hyperref}
\usepackage[backend=biber, sorting=none]{biblatex}
\addbibresource{bibliography.bib}
\usepackage{placeins}
\setlength{\intextsep}{6pt plus 2pt minus 2pt}
\setlength{\abovecaptionskip}{4pt}
\setlength{\belowcaptionskip}{0pt}
% Riduci spazio attorno alle equazioni inline
\setlength{\abovedisplayskip}{4pt}
\setlength{\belowdisplayskip}{4pt}

% Riduci spazio tra paragrafi
\setlength{\parskip}{1pt}

\begin{document}

\begin{center}
{\LARGE Sintesi della Prova Finale}\\[0.3cm]
\small
\textbf{Titolo:} Analisi spettrale di eventi SEP e stime della dose per missioni interplanetarie\\[5pt]
\begin{tabular}{ll@{\qquad}ll}
\textbf{Cognome e Nome:} & Dolci Stefano & \textbf{Matricola:} & 900231 \\
\textbf{Corso di Laurea:} & Fisica & \textbf{Data della seduta:} & 26 Marzo 2026 \\
\textbf{Relatore:} & Prof. Gervasi Massimo & \textbf{Correlatore:} & Dott. Della Torre Stefano \\
\textbf{Recapito telefonico:} & +39 3343715962 & & \\
\end{tabular}\\[0.2cm]
\end{center}
\vspace{0.6cm}
\noindent Nello spazio interplanetario, la radiazione ionizzante è una delle principali minacce per il successo delle missioni spaziali. Al di fuori della magnetosfera terrestre ci sono due le principali sorgenti di radiazione: i raggi cosmici galattici (GCR) e le particelle energetiche solari (SEP). I GCR sono particelle cariche ad alta energia, con origini al di fuori del sistema solare. D'altro canto i SEP sono particelle energetiche, tipicamente protoni, accelerate a seguito di eruzioni solari. La natura imprevedibile dei SEP richiede inevitabilmente un approccio probabilistico quando ci si pone l'obiettivo di stimare il rischio associato. Nel seguente elaborato si stima la dose che un'ipotetica sonda interplanetaria riceverebbe durante una missione verso Marte, considerando il solo contributo dei SEP.\\
Per farlo sono stati selezionati 8 eventi SEP storici dal 1972 al 2024. Per ciascuno di questi sono stati raccolti e processati i flussi differenziali dagli strumenti GOES, IMP-8, ACE e STEREO-A. Integrando questi dati per la durata dell'evento e interpolandoli con una funzione di tipo Band o E-R si ricava lo spettro in energia da \qty{0.05}{\mega\electronvolt} a \qty{1}{\giga\electronvolt}. Le incertezze da propagare sulla dose sono state calcolate con un campionamento Monte Carlo sui parametri del fit. Il codice di trasporto SR-NIEL è stato utilizzato per stimare la \textit{Total Ionizing Dose} (TID) in silicio dietro schermature di alluminio in 2 geometrie diverse: slab planare e guscio sferico, entrambi con spessore equivalente di 10\,g/cm$^2$. 
A ciascun evento si associa la coppia $(J_i, D_i)$, ovvero la fluenza 
integrale omnidirezionale sopra soglia energetica $J(>E_0)$ e la TID calcolata. 
Il fit pesato ai minimi quadrati di questi 8 punti definisce la funzione di trasferimento empirica: $ D(J) = 10^{a} \cdot J^{b}$, valutata con 3 soglie energetiche ($>30,60,100$ MeV). La funzione di trasferimento ottimizzata viene infine applicata alle fluenze cumulative ricavate dal modello probabilistico SAPPHIRE per stimare le dosi attese (1AU) in funzione della durata di missione e fase del ciclo solare.\\
Si stima che per una missione dalla durata di 0.7 anni (259 giorni per orbita Hohmann) durante il massimo solare, con un livello di confidenza del 90\%, la TID in silicio non superi: $37.3\pm3.0$ mGy per la configurazione planare e $62.1\pm6.2$ mGy per la schermatura sferica. Questi valori sono stati confrontati con dei valori medi per il \textit{dose rate} dai soli GCR. A parità di durata della missione e schermatura, durante il massimo solare i SEP dominano la dose con un rapporto SEP/GCR $\sim1.6$, durante il minimo invece i GCR rimangono la componente dominante con un rapporto $\sim0.07$. Questi risultati rappresentano l'\textit{upper bound} conservativo in quanto valutati a distanza fissa 1AU.\\ Con un apposita procedura sviluppata in questo elaborato e simulazioni con il modello di trasporto SOLPENCO2, il risultato è stato riscalato all'orbita di trasferimento Terra-Marte. Considerando quindi gli effetti del trasporto delle particelle lungo l'orbita, è stata stimata una riduzione in percentuale della fluenza, e di conseguenza della dose, pari al $20.5\%$. Tali valori riscalati sono compatibili con le misure \textit{in-situ} dei SEP riportate dalla sonda MSL/RAD durante il transito verso Marte.\\ Il secondo risultato interessante riguarda l'ottimizzazione della funzione di trasferimento $D(J)$. Passare dalla soglia energetica convenzionale $J(>30\,\mathrm{MeV})$, riportata nella maggior parte dei cataloghi, alle soglie ottimizzate $J(>100\,\mathrm{MeV})$ e $J(>60\,\mathrm{MeV})$, riduce lo scatter spettrale da $\pm28-35\%$ a $\pm8-10\%$, creando così una relazione fluenza-dose quasi lineare ($R^2= 0.097$). In questo modo la fluenza integrata condensa tutte le informazioni sulle differenze spettrali e la fisica del trasporto dietro schermatura. Questo permette di ottenere delle stime della dose da SEP ad un costo computazionale inferiore rispetto a simulazioni più complesse. La validità di questa relazione è stata inoltre verificata con \textit{Leave-one-out cross-validation} ($88-100\%$ degli eventi viene predetto entro $\pm1\sigma$ con un RMSE del $10-14\%$) e l'analisi dei residui non mostra correlazione significativa con altri parametri spettrali. \\Le principali limitazioni di queste stime sono: aver considerato solo il contributo dei protoni senza includere particelle secondarie prodotte nella schermatura; il campione degli 8 SEP storici resta statisticamente limitato, aggiungere ulteriori eventi aiuterebbe a controllare la validità della $D(J)$; testare l'effetto di schermature con materiali compositi e spessori diversi. Infine l'uso di codici Monte Carlo più completi (FLUKA/GEANT4) e un modello 3D più complesso di una sonda reale rappresentano il passo successivo naturale di questo elaborato.
\clearpage


\end{document}
